\begin{textsample}{POS Dim 1 – 2021 – Score 23.00 – 2021/t004046.txt}  \label{ex:f1_pos_061}
Bruno Giussani : Hello.
Bill Gates calls himself an imperfect messenger on climate because of his high \textbf{carbon} footprint and lifestyle.
However, he has just made a major contribution to our thinking about confronting climate change via a book, a book about decarbonising our economy and society.
It ' s an optimistic"can-do"kind of book with a strong focus on technological solutions.
He discusses \textbf{the} things we have, such as wind and solar power, \textbf{the} things we need to develop, such as carbon-free cement or carbon-free steel, and long-term energy storage.
And he talks about \textbf{the} economics of it all, introducing \textbf{the} concept of green premium. \textbf{The} gist of \textbf{the} book, and really I ' m simplifying here a lot, is that fighting climate change is going to be hard, but it ' s possible, we can do it.
There is a pathway to a clean and prosperous future for all.
So we want to unpack some of that with \textbf{the} author of \textbf{the} book.
Bill Gates, welcome back to TED.
Bill Gates : Thank you.
Giussani : Bill, I would like to start where you start, from \textbf{the} title of \textbf{the} book."How to Avoid a Climate Disaster,"which, of course, presumes that we are heading towards a climate disaster if we don ' t act differently.
So what is \textbf{the} single most important thing we must do to avoid a climate disaster?
Gates : Well, \textbf{the} \textbf{greenhouse} gases we put into \textbf{the} atmosphere, particularly CO2, stay there for thousands of years.
And so it ' s really \textbf{the} \textbf{sum} of all those \textbf{emissions} are forcing \textbf{the} temperature higher and higher, which will have disastrous effects.
And so we have to take these \textbf{emissions}, which are presently over 51 billion tons per year and drive those all \textbf{the} way down to zero.
And that ' s when \textbf{the} temperature will stop increasing and \textbf{the} disastrous weather events won ' t get worse and worse.
So it ' s pretty demanding.
It ' s not a 50 percent reduction.
It ' s all \textbf{the} way down to zero.
Giussani : Now, 51 billion is a big number.
It ' s difficult to register, to understand.
Help us visualize \textbf{the} scale of \textbf{the} problem.
Gates : Well, \textbf{the} key is to understand all \textbf{the} different sources.
And people are mostly aware of \textbf{the} production of electricity with natural gas and \textbf{coal} as being a big source.
That ' s about 27 percent.
And they ' re somewhat aware of \textbf{transportation}, including passenger cars.
Passenger cars are seven percent and \textbf{transportation} overall is 16 percent.
They have far less awareness of \textbf{the} other three segments.
Agriculture, which is 19 percent, heating and cooling buildings, including using natural gas, are seven percent.
And then, sadly, \textbf{the} biggest segment of all, manufacturing, including steel and cement.
People are least aware of that one.
And in fact, that one is \textbf{the} most difficult for us to solve. \textbf{The} size of all \textbf{the} steel plants, cement plants, paper, plastic, \textbf{the} industrial economy is gigantic.
And we ' re asking that to be changed over in this 30-year period when we don ' t even know how to make that change right now.
Giussani : So your core argument, and really, here I ' m simplifying, is that basically we need to clean up all of that, right? \textbf{The} way we make things, we grow things, we get around and power our economy.
And so to do that, we need to get to a point where green energy is as cheap as \textbf{fossil} fuels, and new materials, clean materials, are as cheap as current materials.
And you call that"eliminating \textbf{the} green premium."So to start, tell us what you mean by green premium.
Gates : Yeah, so \textbf{the} green premium varies from \textbf{emission} sources.
It ' s \textbf{the} cost to buy that product where there ' s been no \textbf{emissions} versus \textbf{the} cost we have today.
And so for an electric car, \textbf{the} green premium is reasonably modest.
You pay a little more upfront, you save a bit on \textbf{the} maintenance and gasoline, you give up some range, you have a longer charging time.
But over \textbf{the} next 15 years, because \textbf{the} volume is there and \textbf{the} R and D is being done, we can expect that \textbf{the} electric car will be preferable.
It won ' t cost more.
It will have a much higher range.
And so that green premium that today is about 15 percent is headed to be zero even without any government subsidies.
And so that ' s magic.
That ' s exactly what we need to do for every other category.
Now, an area like cement, where we haven ' t really gotten started yet, \textbf{the} green premium today is almost double \textbf{the} price.
That is, you pay 125 dollars for a ton of cement today but it would be almost double that if you insisted that it be green cement.
And so \textbf{the} way I think of this is in 2050 we ' ll be talking to India and saying to them,"Please use \textbf{the} green products as you ' re building basic shelter, you know, simple air conditioning,"which they ' ll need because of \textbf{the} heat increase or lighting at night for students.
And unless we ' re willing to subsidize it or \textbf{the} price is very low, they ' ll say,"No, this is a problem that \textbf{the} rich countries created, that India is suffering from and you need to take care of it."So only by bringing that green premium down very dramatically, about 95 percent across all categories, will that conversation go well so that India can make that shift.
And so \textbf{the} key thing here is that \textbf{the} US ' s responsibility is not just to zero out its \textbf{emissions}.
That ' s a very hard thing, but we ' re only 15 percent.
Unless we, through our power of innovation, make it so cheap for all countries to switch all categories, then we simply aren ' t going to get there.
And so \textbf{the} US really has to step up and use all of this innovative capacity every year for \textbf{the} next 30 years.
Giussani : What what needs to happen in order for these breakthroughs to actually \textbf{occur}?
Who are \textbf{the} players who need to come together?
Gates : Well, innovation usually happens at a pace of its own.
Here we have this deadline, 2050.
And so we have to do everything we can to accelerate it.
We need to raise \textbf{the} R and D budgets in these areas.
In 2015 I organized, along with President Hollande and Obama, a side event to \textbf{the} Paris climate talks where what Prime Minister Modi had labeled"mission innovation"was a commitment to double R and D budgets over a five-year period.
And all \textbf{the} big countries came in and made that pledge.
Then we need lots of smart people, who, instead of working on other problems, are encouraged to work on these problems.
So coming up with funding for them is very, very important.
I ' m doing some of that through what I call Breakthrough Energy Fellows.
We need high-risk capital to invest in these companies, even though \textbf{the} risks are very, very high.
And that ' s -- There is now -- Breakthrough Energy Ventures is one group doing that and drawing lots of other people in.
But then \textbf{the} most difficult thing is we actually need markets for these products even when they start out being more expensive.
And that ' s what I call Catalyst, organizing \textbf{the} buying power of consumers and companies and governments so that we get on \textbf{the} learning curve, get \textbf{the} scale going up, like we did with solar and wind across all these categories.
So it ' s both supply of innovation and demand for \textbf{the} green products.
That ' s \textbf{the} combination that can start us to make this change to \textbf{the} infrastructure of \textbf{the} entire physical economy.
Giussani : In terms of funding all this, \textbf{the} financial system as a whole right now is essentially funding expansion of \textbf{fossil} fuels consistent with three degrees Celsius increase in global heating.
What do you say to \textbf{the} Finance Committee, beyond \textbf{the} venture capital, about \textbf{the} need to think and act differently?
Gates : Well, if you look at \textbf{the} interest rates for a solar field versus any other type of investments, it ' s not lower.
You know, money is very fungible.
It ' s going to, you know, different projects.
But, there ' s no, sort of, special rate for climate-related projects.
Now, you know, governments can decide through tax incentives to improve those things.
But you know, this is not just about reporting numbers.
It ' s good to report numbers.
But \textbf{the} steel industry is providing a vital service.
Even gasoline, for 98 percent of \textbf{the} cars being purchased today, allows people to get to their job.
And so, you know, just divestment alone is not going to be \textbf{the} thing that creates \textbf{the} new alternative and brings \textbf{the} cost of that down.
So \textbf{the} finance side will be important because \textbf{the} speed of deployment -- Solar and wind needs to be accelerated dramatically beyond anything we ' ve done today.
In fact, you know, on average, we ' ll have to deploy three times as much every year as \textbf{the} peak year so far.
And so those are a key part of \textbf{the} solution.
They ' re not \textbf{the} whole solution.
But you know, people sometimes think just by putting numbers, disclosing numbers, that somehow it all changes or by divesting that it all changes. \textbf{The} financial sector is important, but without \textbf{the} innovation, there ' s nothing they can do.
Giussani : A lot of \textbf{the} current focus is on cutting \textbf{emissions} by half by 2030 and on \textbf{the} way to reaching that zero by 2050.
And in \textbf{the} book, you write that there is danger in that kind of thinking, that we should keep our main focus on 2050.
Can you explain that?
Gates : Well, \textbf{the} only real measure of how well we ' re doing is \textbf{the} green premium, because that ' s what determines whether India and other developing countries will choose to use \textbf{the} zero \textbf{emission} approach in 2050. \textbf{The} idea, you know -- If you ' re focusing on short-term reductions, then you could say, oh, it ' s \textbf{fantastic}, we just put billions of dollars into natural-gas plants.
And even ignoring that there ' s a lot of leakage that doesn ' t get properly measured, you know, give them full credit and say that ' s a 50 percent reduction. \textbf{The} lifetime of that plant is greater than 30 years.
You financed it, assuming you ' re getting value out of it over a longer period of time.
So to say,"Hey, hallelujah, we switched from \textbf{coal} to natural gas,"that has nothing to do with reaching zero.
It sets you back because of \textbf{the} capital spending involved there.
And so it ' s not a path to zero if you ' re just working on \textbf{the} easy parts of \textbf{the} \textbf{emission}.
That ' s not a path. \textbf{The} path is to take every source of \textbf{emission} and say,"Oh, my goodness, how am I going to get that green premium down?
How am I going to get that green premium down?"Otherwise, getting to 50 percent does not stop \textbf{the} problems.
This is very tough because \textbf{the} nature is zero.
It ' s not just, you know, a small decrease.
Giussani : So getting \textbf{the} green premium down, so you mentioned India.
To make sure that \textbf{the} transition, \textbf{the} clean transition, is also an Indian story and an African story and not only a Western story, should rich countries adopt \textbf{the} expensive clean alternatives now, to kind of buy down \textbf{the} green premium and make technology, therefore, more accessible to low-income countries?
Gates : Absolutely.
You have to pick which product paths with scale will become low green-premium products.
And so you don ' t want to just randomly pick things that are low-emission.
You have to pick things that will come down.
So like, you know, so far, hydrogen fuel cells have not done that.
Now they might in \textbf{the} future.
Solar, wind and lithium-ion, we ' ve seen these incredible cost reductions.
And so we have to duplicate that for other areas, things like offshore wind, heat pump, new ways of doing transmission.
You know, we have to keep \textbf{the} electricity grid reliable, even in very bad weather conditions.
And so that ' s where storage or nuclear and transmission will have to be scaled up in a very significant way, which we now have this \textbf{open-source} model to look at that.
And so \textbf{the} buying of green products that demand signal, what I call Catalyst, we ' re going to have to orchestrate a lot of money, many tens of billions of dollars for that, that is one of \textbf{the} most expensive pieces so that \textbf{the} improvements come in these other areas and some of them we haven ' t even gotten started on.
Giussani : Many people believe actually that \textbf{the} climate question is mostly a question of consuming less and particularly consuming less energy.
And in \textbf{the} book, you actually write several times that we need to consume more energy.
Why?
Gates : Well, \textbf{the} basic living conditions that we take for granted should be made available to all humans.
And \textbf{the} human population is growing.
And so you ' re going -- as you provide shelter, heating and air conditioning, which anywhere near \textbf{the} equator will be more in demand since you ' ll have many days where you can ' t go outdoors at all.
So we ' re not going to stop making shelter.
We ' re not going to stop making food.
And so we need to be able to multiply those processes by zero.
That is, you make shelter, but there ' s no \textbf{emissions}.
You make food, but there ' s no \textbf{emissions}.
That ' s how you get all \textbf{the} way down to zero.
Now, it ' s made somewhat easier if rich countries are consuming less, but how far will that go?
Giussani : So if I summarize in my head your book, you are basically suggesting that if we eliminate \textbf{the} green premium, \textbf{the} transition somehow can \textbf{occur}, \textbf{the} cost and technology, green premium and breakthroughs are \textbf{the} key drivers.
But, you know, you think of climate and climate is kind of a wicked problem.
It has implications that are social, political, behavioral.
It requires significant citizen involvement.
Are you focusing too much on tech and not enough on those other variables?
Gates : Well, we need all these things.
If you don ' t have a deep engagement, particularly by \textbf{the} younger generation making this a top priority every year for \textbf{the} next 30 years across all \textbf{the} developed countries, we will not succeed.
And I ' m not \textbf{the} one who knows how to activate all those people.
I ' m super glad that \textbf{the} people who are smart about that are thinking.
It ' s a necessary element.
Likewise, \textbf{the} piece that I do have experience in, innovation ecosystems, that ' s a necessary element.
You will not get there just by saying,"Please stop using steel."You know, it won ' t happen.
And so \textbf{the} innovation piece has to come along and particularly encouraging consumers to buy electric cars or \textbf{artificial} meat or electric heat pumps, they ' re part of driving that demand, what I call \textbf{the} catalytic demand.
That alone won ' t do it, because some of these big projects, like green hydrogen or green aviation fuel require billions in capital expense.
But \textbf{the} demand signal from enlightened consumers is very important.
Their political voice is very important.
Their pushing \textbf{the} companies they work at is very important.
And so I absolutely agree that \textbf{the} broad community who cares about this, particularly if they understand how hard it is and don ' t say,"Oh, we can do it in 10 years,"they are super important.
Giussani : You mentioned \textbf{artificial} meat and I know you love burgers.
Have you tried an \textbf{artificial} burger, and how did you find it?
Gates : Yes, I ' m an investor in all these Impossible, Beyond and various people.
And I have to say, \textbf{the} progress in that sector is greater than I expected.
Five years ago, I would have said that is as hard as manufacturing.
Now it ' s very hard, but not as hard as manufacturing.
There is no Impossible Foods of green steel and \textbf{the} quality is going to keep improving.
It ' s quite good today.
You know, there ' s other companies coming in to that space covering different types of food.
And as \textbf{the} volume goes up, \textbf{the} price will go down.
And so I think it ' s quite promising.
I have to admit, 100 percent of my burgers aren ' t \textbf{artificial} yet, it ' s about 50 percent, but I ' ll get there.
Giussani : It ' s a start.
You ' re already on 2030 on that on that front.
So you mentioned your expertise in innovation.
You have a couple of lines of -- how to say it? -- of exquisite modesty in \textbf{the} book where you say you think like an engineer, you don ' t know much about politics, but actually you talk to top politicians more than any of us.
What do you ask them and what do you hear in return?
Gates : Well, they ' re mostly responding to voters interest.
Most of \textbf{the} countries we ' re talking about are democracies.
And you know, there are resources that will need to be put in like tax credits that encourage buying green products or government purchasing of green products at slightly higher prices.
Those are real trade-offs.
And, you know, making sure that \textbf{the} areas where there ' s lots of jobs in, say, \textbf{coal} mining, or things where \textbf{the} demand will go down, having that political sensitivity and thinking through, can we put \textbf{the} new jobs in that area and what are \textbf{the} departments that handle that.
This is a tough political problem.
And, you know, my admonition to people is not only to get educated themselves, but help educate other people.
And often, like in \textbf{the} US, if it ' s people of both parties, that ' s even better. \textbf{The} level of interest is high, but it needs to get even higher, almost like a moral mission of all young people to go beyond their individual success, that they believe that getting to zero by 2050 is critical.
Giussani : Thank you.
Now, before we continue \textbf{the} interview, I would like to take a short detour for one minute because my colleagues at TED-Ed have produced a series of seven animated videos inspired by your book.
They introduce \textbf{the} concept of net zero \textbf{emissions}, they discuss other questions and challenges that we are facing relating to climate.
So I would like to share one short \textbf{clip} from one of those animations.
You flip a switch, \textbf{coal} burns in a furnace which turns water into steam.
That steam spins a \textbf{turbine} which activates a generator which pushes electrons through \textbf{the} wire.
This current propagates through hundreds of miles of electric cables and arrives at your home.
All around \textbf{the} world, countless people are doing this every second : flipping a switch, plugging in, pressing an"on"button.
So how much electricity does humanity need?
Giussani : \textbf{The} answer to that and to many other climate questions, of course, is in \textbf{the} seven animations available on \textbf{the} TED-Ed site, and \textbf{the} TED-Ed YouTube channel.
Bill, you mentioned \textbf{the} younger generation before.
How does \textbf{the} younger generation ' s role in solving this problem inspire you?
Gates : Well, they... can make sure this is a priority.
And if we have, you know, it ' s a priority for four years, then it ' s not for four years, you can ' t ask \textbf{the} trillions of investment in \textbf{the} new approach to take place.
It ' s got to be pretty clear that even though political parties may disagree on \textbf{the} tactics, \textbf{the} same way they agree there should be a strong defense that they agree this zero by 2050 is a shared goal, and then discuss, OK, where does government come in, where does \textbf{the} private sector come in, which area deserves priority?
That will be a huge milestone where it ' s a discussion about how to get there versus whether to get there.
And \textbf{the} US is \textbf{the} most fraught in terms of it being politicized, even in terms of, is this a huge problem or not.
So, you know, young people, they ' re going to be around to see \textbf{the} good news if we ' re able to achieve this.
You know, I won ' t be around.
But, they speak with moral authority and they have particular people who are stepping up on this.
But I hope that ' s just \textbf{the} beginning.
And that ' s why this year, I think, is so important.
All this recovery money being programed, people thinking about, do governments protect us \textbf{the} way they should, do governments work together?
You know, \textbf{the} pandemic has teed this up, it ' s OK, what ' s \textbf{the} next big problem we need to collaborate around?
And I ' m hoping that climate appears there because of these activists.
Giussani : Yeah, we ought to come back to \textbf{the} pandemic question.
But I want to talk a moment about some specific technology breakthroughs that you mention in \textbf{the} book.
But you also just mentioned \textbf{the} organization you set up, Breakthrough Energy, to invest in clean tech start-up and advocate for policies.
And I assume somehow your book is kind of a blueprint for what Breakthrough Energy is going to do.
There is a branch in that organization, called Catalyst, that ' s prioritizing several technologies, including green hydrogen direct carbon-capture, aviation biofuels.
I don ' t want to ask you about specific investments, but I would like to ask you to describe why those priorities, maybe starting with green hydrogen.
Gates : If we can get green hydrogen that ' s very cheap, and we don ' t know that we can, that becomes a magic ingredient to a lot of processes that lets you make fertilizer without using natural gas, it lets you reduce \textbf{iron} ore for steel production without using any form of \textbf{coal}.
And so there ' s two ways to make it.
You can take water and split it into hydrogen and oxygen.
You can take natural gas and pull out \textbf{the} hydrogen.
And so it ' s kind of a holy grail, you know, and so we need to get going.
We need to get all \textbf{the} components to be very, very cheap.
And only by actually doing projects, significant scale projects, do you get on to that learning curve.
And so Catalyst will fund \textbf{the} early pilot projects along with governments to go make green hydrogen, get \textbf{the} electrolyzers to get up in volume and get a lot cheaper, because that would be a huge advance.
Lot of manufacturing, not all of it, but a lot of it would be solved with that.
Giussani : So just for those who don ' t know, green hydrogen is produced with clean energy sources, wind, solar and so, nitrogen produced from natural gas or \textbf{fossil} gas is called gray hydrogen. \textbf{The} second priority that you have set is direct \textbf{carbon} capture, pulling \textbf{carbon} from \textbf{the} atmosphere.
And at least in theory, we can find a way to do that at large scale.
Together with other technologies \textbf{the} problem could be solved, but it ' s a very early-stage, unproven technology.
You describe it yourself in \textbf{the} book as a thought experiment at this stage.
But you are one of \textbf{the} main investors in this sector in \textbf{the} world.
So what ' s \textbf{the} real potential for direct \textbf{carbon} capture?
Gates : So there ' s a company today, Climeworks, that for a bit over 600 dollars a ton will do capture.
Now it ' s at fairly small scale.
I ' m a customer of theirs as part of my program where I eliminate all of my \textbf{carbon} \textbf{emissions} in a gold standard way.
There are other people who are trying to do larger-scale plants like Carbon Engineering.
Breakthrough Energy is investing in a number of these carbon-capture entities.
In a way, \textbf{the} \textbf{carbon} capture is for \textbf{the} part that you can ' t solve any other way.
It ' s kind of \textbf{the} brute-force piece, and no one knows what that price will be.
If it ' s 100 dollars a ton, then \textbf{the} cost against current \textbf{emissions} would be five trillion a year.
Can we get below 100 dollars a ton?
It ' s not clear that we can, but it would be \textbf{fantastic} to take \textbf{the} final 10, 15 percent of \textbf{emissions}, and instead of making \textbf{the} change at \textbf{the} place where \textbf{the} \textbf{emission} takes place, just do this, direct air capture.
To be clear, direct air capture means you ' re just filtering \textbf{the} air, and you ' re pulling out \textbf{the} 410 million current parts per million and putting that into a pressurized form of CO2 that then you sequester in someplace that you know it ' ll stay for millions of years.
And so this industry is just at \textbf{the} very, very beginning.
But it ' s a necessary piece for \textbf{the} tail of \textbf{emissions}.
Giussani : And a third priority is aviation biofuels.
Now flying over \textbf{the} last several years has become a symbol of a \textbf{polluting} lifestyle, let ' s say.
Why aviation biofuels as a priority?
Gates : Well, \textbf{the} great thing about passenger cars is that even though batteries don ' t store energy as well as gasoline does, there ' s dramatic difference, you can afford to have \textbf{the} extra size and weight of those batteries.
On a plane that is a large plane going a long distance, there ' s no chance that batteries will ever have that energy density.
I mean, there ' s some crazy people who are working on it, and I ' ll be glad to fund them, but you wouldn ' t want to count on that.
And so you either have to use hydrogen as a plane fuel, that has certain challenges, or you just want to make today ' s aviation fuel with green processes, like using plants as \textbf{the} source of how you make those.
And so I ' m \textbf{the} biggest individual customer of a group that makes green aviation fuel.
So that ' s what I ' m using now, costs over twice as much as \textbf{the} normal aviation fuel.
But as \textbf{the} demand scales up, that ' s one of those green premiums that we hope to bring down so that at some point lots of consumers will say,"Yes, I ' ll pay a little extra on my plane \textbf{ticket},"probably through Catalyst or through \textbf{the} airline to make sure that we ' re building up that industry and trying to get \textbf{the} costs, that green premium down quite a bit.
So this is a super important area that, you know, I was amazed when I went to buy that I was by far \textbf{the} biggest individual customer.
Giussani : It ' s also, of course, a very symbolic area, right?
People think of cars and think of \textbf{airplanes} when they think of pollution and \textbf{emissions}.
There is a fourth technology that I want to bring up because you are an advocate of and an investor in nuclear power, which is also not universally accepted as clean energy because of \textbf{the} risks, because of radioactive waste.
Can you make \textbf{the} case, briefly, but can you make your case for nuclear?
Gates : Well, one thing that needs to be appreciated is that energy has to come from somewhere.
And so as you stop using natural gas to heat homes and gasoline to power cars, \textbf{the} electric grid will have to grow dramatically.
And even in \textbf{the} case of \textbf{the} US, where electricity demand has been flat \textbf{the} last few decades, will need to be almost three times as large.
As you do that with weather-dependent sources, \textbf{the} reliability of \textbf{the} electricity generation goes away, that as you get big cold fronts that for 10 days can stop most of \textbf{the} wind and solar, say, in \textbf{the} Midwest.
So \textbf{the} question is, how do you use massive amounts of transmission and storage and non weather-dependent sources which at scale nuclear as \textbf{the} only choice there, to maintain that reliability and not have people, say, freeze to death.
And nuclear -- People will be pretty impressed with how valuable it is to create that reliability.
Now, 80 percent, at least in \textbf{the} US, will be wind and solar.
So we have to build that faster than ever.
But what is that piece that ' s always available is where nuclear would fit in.
Giussani : Now, you talk a lot about scaling up technologies in \textbf{the} book.
I want to ask you a question about scaling down, because one thing that ' s absent from your book is a discussion about scaling down \textbf{fossil} fuel production, maybe starting with at least not expanding it anymore, with adopting a sort of \textbf{fossil} fuel nonproliferation approach, because right now we are talking about a green transition, but at \textbf{the} same time we are building out further \textbf{fossil} fuel infrastructure.
What ' s your view on that?
Gates : Well, it ' s all about \textbf{the} demand for \textbf{fossil} fuels and \textbf{the} fact that \textbf{the} green premium is very high.
If you want to restrict supply, then you ' ll just drive \textbf{the} price up.
People still want to drive to their job.
If you know, say you made \textbf{fossil} fuels illegal.
You know, try that for a few weeks.
My view is you have to create a substitute because \textbf{the} services provided by \textbf{the} \textbf{fossil} fuel are actually quite valuable.
Electricity is valuable.
Transportation is valuable.
And are you willing to drop demand for those things to zero?
So \textbf{the} capitalistic economy will respond.
If it ' s clear, you know, how you ' re going to do your tax policies and your credit, you ' re going to drive innovation, then \textbf{the} infrastructure investments in those things will go down.
We see that with \textbf{coal} already today.
But...
Just speaking against something when you haven ' t created an alternative isn ' t going to get you to zero.
Giussani : OK, I want to shift topics.
You mentioned before \textbf{the} pandemic, and your main focus since you set up \textbf{the} Gates Foundation has been on global health.
I actually read \textbf{the} letter, \textbf{the} annual letter you wrote with your wife, Melinda, in January 2021, which was very much about \textbf{the} COVID-19 pandemic.
And I found myself marveling at how what you write resonates with \textbf{the} climate challenge.
What lessons from \textbf{the} pandemic can be actually applied to climate?
Gates : Well I think \textbf{the} biggest lesson is that governments got to, on our behalf, avoid disastrous future outcomes.
And individual citizens aren ' t equipped to either do those evaluations or think through that R and D and deployment plan that ' s necessary here.
And so we have to make it an imperative that governments, hopefully of any party, join in to this. \textbf{The} pandemic, we did eventually get global cooperation. \textbf{The} US didn ' t play its normal role there, but \textbf{the} private sector innovation created \textbf{the} vaccine.
Now, sadly, with climate, \textbf{the} pain it ' s causing gets worse over time.
So with \textbf{the} pandemic, we had all these deaths, and people were like,"Wow, OK, we should do something."With climate you can ' t wait.
You know, \textbf{the} coral reefs will have died off, \textbf{the} species will be gone.
And so if you say,"OK, well, let ' s, when it gets bad, we ' ll invent something like a vaccine."That doesn ' t work, because to stop \textbf{the} \textbf{emissions}, you have to change every steel plant, cement plant, car, things that have massive lead times and literally trillions of dollars of investment.
And so with \textbf{the} pandemic, we messed up.
We didn ' t pay attention.
You know, people like myself said it was a problem.
But, now we ' re getting our way out of it through innovation.
But climate ' s harder, much harder problem.
And so \textbf{the} political will to get it right needs to be unprecedented compared to \textbf{the} pandemic or almost any other political cause.
Giussani : I want to correct you about private-sector innovation, because, of course, a lot of public money went into funding research and guaranteed purchases.
And so for \textbf{the} vaccines.
But that ' s a separate interview.
Gates : Well, \textbf{the} Pfizer vaccine used no government money.
Giussani : That ' s absolutely true.
Bill, still related to our reaction to these big challenges, over half of \textbf{the} total \textbf{greenhouse} gas \textbf{emissions} since 2017, 50 or 51 have gone up in \textbf{the} atmosphere in \textbf{the} last 30 years.
And we have known for more than 30 years that they have pernicious effects, to say \textbf{the} least.
So if we could mobilize such enormous resources and policies and collaborations and global collaborations against COVID-19, what ' s \textbf{the} lever we can use to mobilize \textbf{the} same against climate?
Gates : Well, until 2015, nobody pushed for \textbf{the} R and D budgets to go up.
So, you know, there are times when I think"Wow, are we serious, are we not?"You know, that R and D question just wasn ' t discussed.
And you could read every paper on green steel and there were hardly any.
And so we are just starting to get serious about climate.
And you know, we ' ve wasted a lot of time that we should have used to work on \textbf{the} hard parts of climates, just having short-term goals and not focusing on \textbf{the} R and D piece, you know, \textbf{the} last 20 years, we don ' t have much to show for \textbf{the} hard categories.
Now, there ' s still time to get there.
We will have to fund adaptation a great deal because \textbf{the} 2050 goal is not because that ' s a goal that gives you zero damage.
It ' s simply \textbf{the} goal that is \textbf{the} most ambitious that has a chance of being achieved.
And we are causing problems for subsistence farmers and sea-level rise and wildfires and natural ecosystems.
And so \textbf{the} adaptation side is even more underfunded than \textbf{the} mitigation side.
For example, helping poor farmers with seeds that can deal with \textbf{the} \textbf{droughts} and high temperatures is funded at less than a billion a year, which is deeply tragic.
Giussani : Bill, we ' re getting towards \textbf{the} end.
I mentioned at \textbf{the} beginning, in \textbf{the} book you describe yourself as an imperfect messenger on climate.
What changes have you made in your personal and family life to reduce your footprint?
Gates : Well, I ' m certainly driving an electric car, you know, putting solar panels on \textbf{the} houses where that makes sense.
Using this green aviation fuel, I still can ' t say that I ' ve stopped eating meat or that I never fly.
And so it ' s mostly by funding at over seven million a year \textbf{the} products that although their green premiums are very high today, like \textbf{the} \textbf{carbon} capture or like \textbf{the} aviation fuel, by funding those, you actually get those on to \textbf{the} learning curve to get those prices down very dramatically.
One area that I do is I fund putting in electric heat-pumps into low-cost housing.
And so instead of using natural gas, they get a lower bill because it ' s done with electricity.
And I paid that extra capital cost as an offset.
So, you know, accelerating those demand things.
I ' ve tried to get out in front of that.
Giussani : You also talk in \textbf{the} book about paying more than market value or \textbf{the} current market price for offsets, and you hinted to them before, talking about \textbf{the} golden standard.
Tell us about what kind of -- offsets can be kind of confusing and controversial.
What are \textbf{the} right offsets?
Gates : Well, first of all, it ' s great that we ' re finally talking about offsets, and any company that actually looks at their \textbf{emissions} and pays for offsets is way better than a company that either doesn ' t look at their \textbf{emissions} or looks at them and doesn ' t pay for offsets.
And so \textbf{the} leading companies are now buying offsets, some of them spending hundreds of millions to buy offsets.
That is a very, very good thing. \textbf{The} ability to see which offsets actually have long-term benefit that really keep \textbf{the} \textbf{carbon} out for \textbf{the} over thousands of years that count.
There are now organizations that I and others are funding to label offsets as either gold standard or different levels of impact.
And \textbf{the} price of offsets range from 15 dollars a ton to 600 dollars a ton.
Some of \textbf{the} low-cost ones may be really legitimate, like reducing natural gas leakage.
That is pretty dramatic in some cases in terms of \textbf{the} dollars per tons avoided there.
A lot of \textbf{the} forestry things will probably not end up looking that good, as you really look at \textbf{the} lifetime of \textbf{the} tree or what would have happened otherwise.
But, you know, at least we ' re talking about offsets now.
And now we ' re going to do it in a thoughtful way.
Giussani : So some of \textbf{the} people listening to this interview may already be very involved with climate, others may be looking for ways to step up.
And at \textbf{the} end of your book, you have a chapter about individual action.
Give us maybe, say, two examples, two practical examples of things that individual citizens in \textbf{the} US, but also elsewhere, can do to play a meaningful role in tackling climate change.
Gates : I think everybody should start by learning more, you know.
How much steel do we make and where are those steel plants? \textbf{The} industrial economy is kind of a miracle, although sadly, it ' s a source of so many \textbf{emissions}.
Once you really educate yourself, then you ' re in a position to educate others, hopefully, of diverse political beliefs about why this is so important.
And yet it ' s also very, very hard to do.
You have all your buying behavior, electric cars, \textbf{artificial} meat, you know, and you ' ll see for all \textbf{the} different products various things that indicate how green that product is.
And your demand doesn ' t just save those \textbf{emissions}, it also encourages \textbf{the} improvement of \textbf{the} green product.
Political voice, I ' d still put it number one.
Making sure your company is measuring its \textbf{emissions} and is starting to fund offsets and is willing to be a customer for breakthrough storage solutions or green aviation fuel.
That is catalytic.
And a lot of those funds hopefully will go through this vehicle that ' s really identifying which projects globally are technologies that won ' t stay expensive but can do like what solar and wind did and come down in price.
So individuals are what drive this thing. \textbf{The} democracies are where most of \textbf{the} innovation power is and they have to get activated and set \textbf{the} example.
Giussani : I would like to end on \textbf{the} outcome.
Maybe we should have started with \textbf{the} outcome, but let ' s imagine a world where we will actually have done all of what you describe in \textbf{the} book and everything else that ' s necessary.
What would that future everyday life look like?
Gates : Well, I think everybody will be really proud that humanity came together on a global basis to make this radical change and there ' s no precedent for it.
Even world wars, where we orchestrated lots of resources, it was, like, four or five year duration.
And here we ' re talking about three decades of hard work dealing with an enemy that \textbf{the} super bad stuff is out in \textbf{the} future.
And so you ' re benefiting young people and future generations.
In some ways, you know, life will look a lot like it does today.
You ' ll still have buildings, you ' ll have air conditioning, you ' ll have lights at night.
But all of those you ' ll multiply by zero in terms of what \textbf{the} \textbf{emissions} that come out of those activities are.
During \textbf{the} same time frame we ' ll have advances in medicine of curing cancer and finishing polio and \textbf{malaria} and all sorts of things.
And so by taking away this one super negative thing, then all \textbf{the} progress we make in other areas won ' t get reduced by this awful thing that, if it goes unchecked, \textbf{the} migration out of \textbf{the} equatorial regions, \textbf{the} number of deaths, it ' ll make \textbf{the} pandemic look like nothing.
I think so from a moral point of view and letting \textbf{the} other improvements not be offset by this.
It ' ll be a source of great pride that, hey, we came together.
Giussani : OK, let ' s hope that we actually do.
Bill Gates, thank you for sharing your knowledge, thank you for this very, very important book.
And thank you for coming back to TED.
I want to close by showing another short \textbf{clip} from \textbf{the} TED-Ed animation series inspired by Bill ' s book.
This one is about material that ' s all around us, has a big \textbf{carbon} footprint and how to reinvent it.
Now, \textbf{the} series includes seven videos.
You can watch them all for free at ed. ted. com / planforzero.
That ' s planforzero, one single word.
And on \textbf{the} TED-Ed YouTube channel.
Thank you.
Goodbye.
Look around your home.
Refrigeration, along with other heating and cooling, makes up about six percent of total \textbf{emissions}.
Agriculture, which produces our food, accounts for 18 percent.
Electricity is responsible for 27 percent.
Walk outside and \textbf{the} cars zipping past, planes overhead, trains ferrying commuters to work.
Transportation, including \textbf{shipping}, contributes 16 percent of \textbf{greenhouse} gas \textbf{emissions}.
Even before we use any of these things, making them produces \textbf{emissions}, a lot of \textbf{emissions}.
Making materials, concrete, steel, plastic, glass, aluminum and everything else accounts for 31 percent of \textbf{greenhouse} gas \textbf{emissions}. [ Watch \textbf{the} full animated series at ed. ted. com / planforzero and youtube. com / ted-ed ]

% matched lemmas: airplane, artificial, carbon, clip, coal, drought, emission, fantastic, fossil, greenhouse, iron, malaria, occur, open-source, pollute, shipping, sum, the, ticket, transportation, turbine
\end{textsample}
